\documentclass[11pt]{article}
\usepackage[margin=1in]{geometry}
\usepackage[T1]{fontenc}
\usepackage[utf8]{inputenc}
\usepackage{amsmath,amssymb,amsthm}
\usepackage{enumitem}

\title{ICPC World Finals 2008\\G. Net Loss}
\author{}
\date{}

\begin{document}
\maketitle

\section*{Problem Summary}

State the problem in your own words. Focus on the mathematical or algorithmic core rather than repeating the full statement.

\section*{Key Observations}

\begin{itemize}[leftmargin=*]
    \item Write the structural observations that make the problem tractable.
    \item State any useful invariant, monotonicity property, graph interpretation, or combinatorial reformulation.
    \item If the constraints matter, explain exactly which part of the solution they enable.
\end{itemize}

\section*{Algorithm}

\begin{enumerate}[leftmargin=*]
    \item Describe the data structures and the state maintained by the algorithm.
    \item Explain the processing order and why it is sufficient.
    \item Mention corner cases explicitly if they affect the implementation.
\end{enumerate}

\section*{Correctness Proof}

We prove that the algorithm returns the correct answer.

\paragraph{Lemma 1.}
State the first key claim.

\paragraph{Proof.}
Provide a concise proof.

\paragraph{Lemma 2.}
State the next claim if needed.

\paragraph{Proof.}
Provide a concise proof.

\paragraph{Theorem.}
The algorithm outputs the correct answer for every valid input.

\paragraph{Proof.}
Combine the lemmas and finish the argument.

\section*{Complexity Analysis}

State the running time and memory usage in terms of the input size.

\section*{Implementation Notes}

\begin{itemize}[leftmargin=*]
    \item Mention any non-obvious implementation detail that is easy to get wrong.
    \item Mention numeric limits, indexing conventions, or tie-breaking rules if relevant.
\end{itemize}

\end{document}
